%!TEX root = TQFTbook.tex
%%%%%%%%%%%%%%%%%%%%%%%%%%%%%%
\chapter{Algebraic prerequisites: monoidal categories}\label{c:categories}

In this chapter we introduce some of the algebraic language
necessary for describing TQFTs, most importantly the notion
of a monoidal category. We only give a very brief introduction,
referring the reader to the books \ocite{mac-lane} and \ocite{etingof-fusion}
for details.

%%%%%%%%%%%%%%%%%%
\section{Category theory basics}\label{s:category-basics}
We assume that the reader is familiar with basic notions of categories and
functors. We will completely ignore all set-theoretic difficulties related
to the definition of a category; if necessary, the reader can assume that
all sets of objects, morphisms, etc., in all the categories discussed are
subsets of a large enough universal set $U$.

For a category $\calC$, we will use the notation $X\in \calC$ to show that
$X$ is an object of $\calC$. For two objects $X,Y\in\calC$, we denote by
$\Hom_\calC(X,Y)$ the set of morphisms from $X$ to $Y$ (we will drop
the subscript $\calC$ when there is no ambiguity).

Given categories $\calC, \calD$, we denote by $\Fun(\calC, \calD)$ the
category of all functors $\calC\to \calD$.

Examples:
\begin{itemize}
  \item $\sets$: category of all sets.
  \item $\smoothman_n$: category of smooth manifolds of dimension $n$.
    There are different versions: oriented manifolds, framed manifolds, etc.
  \item $\plman_n$: category of piecewise-linear (PL) manifolds of
    dimension $n$.
     Note that
    in dimensions $n\le 6$, every PL manifold has a canonical smooth structure
    and vice versa. In general, it is not so: every smooth manifold has
    a unique PL structure (Whitehead's theorem), but a PL manifold might
    not have a compatible smooth structure, or if it has one, it might not
    be unique (Milnor's exotic spheres). We will discuss PL manifolds in
    more detail later (see \chref{c:PL}).
  \item $\vect_\kk$: category of vector spaces over a field $\kk$.
  \item $R\mod$: category of left modules over an associative ring $R$.
\end{itemize}

One special case of a category, which we will use frequently, is a groupoid.
\begin{definition}\index{groupoid}
  A groupoid is a category in which all morphisms are invertible.
\end{definition}
\begin{example}\label{x:fundamental-groupoid}
  Let $X$ be a topological space. Define its \emph{fundamental
  groupoid}\index{fundamental groupoid} $\Pi_1(X)$ as the category whose objects
  are points in $X$ and morphisms are homotopy classes (relative to the
  endpoints) of paths between points, with the obvious composition.
\end{example}

One common theme in all our constructions is equality versus isomorphism. When
talking about points in a set (e.g., elements of a ring), we can require them to
be equal; for example, we can require $a+b=b+a$ for elements of a ring. However,
while it is technically possible to require two objects in a category to be
equal, it rarely makes sense to do so. For example, in the category of finite
sets, sets $A\times B$ and $B\times A$ are not equal. On the other hand,
requiring two objects to be isomorphic is typically too weak: e.g., any two sets
with the same number of elements are isomorphic.

The common way to resolve this is the following convention.
\begin{convention}\label{conv:canonical}
  A collection of objects $X_\al$ of a category $\calC$, together with a
  collection of isomorphisms $f_{\be\al}\colon X_\al \isoto X_\be$, such that
  $f_{\al\be}f_{\be\ga}=f_{\al\ga}$, $f_{\al\al}=\id_{X_\al}$, is treated as a
  single object in $\calC$.
\end{convention}
This convention, though rarely stated explicitly, is actually widespread.
Virtually every math textbook treats sets $A\times B$ and $B\times A$ as the
same set, even though they are technically not the same object of the category
$\sets$; from the point of view of the convention above, it is justified by
noting that we have a canonical isomorphism $A\times B\isoto B\times A\colon
(a,b)\mapsto (b,a)$.

Formally, this convention can be justified by constructing a category
$\widehat\calC$, whose objects are such collections $(X_\al, f_{\al\be})$, and
observing that the categories $\calC$ and $\widehat\calC$ are equivalent; see,
e.g., \ocite{bakalov-kirillov}*{Lemma 1.1.12}.


%%%%%%%%%%%%%%
\section{Abelian categories}\label{s:abelian}
\begin{definition}\label{d:abelian}\index{abelian category}
  An abelian category is a category in which for every pair of objects $X,Y$,
  the set of morphisms $\Hom(X,Y)$ is an abelian group, composition is
  additive, we have a zero object, the direct sum of two objects, and for every
  morphism, we have the notion of kernel, image, and cokernel satisfying the
  usual properties (see \ocite{etingof-fusion}*{Section~1.3} for details).
\end{definition}
In such a category, we can define the notion of a simple object (a nonzero
object which has no subobjects other than $0$ and itself), a semisimple one
(isomorphic to a direct sum of simple objects), etc. We can also define the
notions of short exact sequences, Jordan--H\"older series (which in turn allows
one to define the notion of the length of an object), injective and projective
objects, and all other usual constructions of homological algebra.

An abelian category is
\emph{$\kk$-linear}\index{category!kk-linear@$\kk$-linear} if all sets
$\Hom(X,Y)$ are vector spaces over a field $\kk$, and composition of morphisms
is bilinear over $\kk$.

Unless explicitly stated otherwise, functors between abelian categories are
assumed to be additive, and functors between $\kk$-linear categories are in
addition assumed to be $\kk$-linear on morphisms.


\begin{example}\label{x:linear-abelian-cats}
  Below are some examples of $\kk$-linear abelian categories. In all cases,
  the morphisms are the obvious ones.
  \begin{enumerate}
    \item $\vect_\kk$;
    \item category $\Rep_\kk(G)$ of representations of a group $G$;
    \item category $\Rep_\kk(\g)$ of representations of a Lie algebra $\g$;
    \item category $A\mod$ of modules over an associative algebra $A$
      over $\kk$.
  \end{enumerate}
\end{example}

Typically, one imposes some finiteness conditions on $\kk$-linear
categories. We will usually assume that all linear categories, unless
explicitly stated otherwise, are
\emph{locally finite}\index{category!locally finite}, i.e.,
\begin{itemize}
  \item All spaces $\Hom(X,Y)$ are finite-dimensional
  \item Every object has finite length
\end{itemize}
(see \ocite{etingof-fusion}*{Section 1.8}).

One can also impose stricter requirements, for example semisimplicity (every
short exact sequence splits), finiteness (there are finitely many isomorphism
classes of simple objects), and more. We will discuss these in more detail when
they become necessary.

One useful fact about $\kk$-linear categories is the following version of the
famous Schur's lemma in representation theory.

\begin{lemma}\label{l:schur}
  Let $\calA$ be a locally finite $\kk$-linear category.
  \begin{enumerate}
    \item If $X,Y\in \calA$ are simple objects, then any morphism
      $f\colon X\to Y$ is either zero or invertible.
    \item If $\kk$ is algebraically closed, then for any non-zero
      simple $X\in \calA$, $\Hom_\calA(X,X)=\kk$.
  \end{enumerate}
\end{lemma}
As with the usual Schur's lemma, part 1 is trivial, and part 2 is immediate
from the fact that the only finite-dimensional division algebra over an
algebraically closed field $\kk$ is $\kk$ itself (see, e.g.,
\ocite{pierce}*{Lemma 3.5}).

%%%%%%%%%%%%%%%%%%%%%%%%%%%%%%%
\section{Monoidal categories}\label{s:monoidal}
One of the most important notions is the notion of a monoidal category, which
informally is described as a category with an associative operation; the model
example is the category of vector spaces with the operation of tensor product.
In this section, we review this notion; as before, we refer the reader to
\ocite{etingof-fusion}*{Chapter 2} and \ocite{leinster}*{Section 1.2} for a more
in-depth discussion.

Before defining monoidal categories, let us look at a simpler notion: a monoid.

A \emph{monoid}\index{monoid} is a set together with an associative
multiplication operation and a unit:
\begin{align*}
  &(a\cdot b)\cdot c = a\cdot (b\cdot c)\\
  &1\cdot a = a\cdot 1 = a
\end{align*}
(if we also add the requirement that every element has an inverse,
we get the definition of a group).

\begin{theorem}\label{t:coherence-monoids}
  Let $M$ be a monoid and $a_1, \dots, a_n\in M$. Consider all possible
  expressions obtained by writing $a_1, \dots, a_n$ \textup{(}in this
  order\textup{)}, then inserting some copies of $1$ and then putting
  parentheses so that multiplications are just binary products, e.g.,
  $(a_1 \cdot 1) \cdot (1 \cdot a_2) $. Then all of them are equal: if
  $X_1, X_2$ are two such expressions \textup{(}with the same
  $a_1, \dots, a_n$\textup{)}, then $X_1=X_2$.
\end{theorem}
It may seem obvious, but in fact it is not quite trivial; essentially it is the
statement that if we construct a graph whose vertices correspond to different
expressions written above, and edges to applications of the associativity and
unit relations, then this graph is connected.

If we only include $a_1, \dots, a_n$ but no 1's, then such expressions are
naturally described by planar binary trees. For example, the bracketings $((a_1
a_2)a_3) a_4$ and $(a_1 a_2)(a_3 a_4)$ correspond to the trees below:

\begin{center}
  %%%%% ((12)3) 4
  \begin{tikzpicture}
    %top level
    \node (1) at (0,0){};
    \node (2) at (8mm,0){};
    \node (3) at (16mm,0){};
    \node (4) at (24mm,0){};
    % labels
    \node[above, draw = none] at (1) {$a_1$};
    \node[above, draw = none] at (2) {$a_2$};
    \node[above, draw = none] at (3) {$a_3$};
    \node[above, draw = none] at (4) {$a_4$};
    % 2nd level
    \node[dotnode] (12) at (4mm, -6mm){};
    \draw (1)--(12)--(2);
    %3rd level
    \node[dotnode] (123) at (8mm, -12mm){};
    \draw (12)--(123)--(3);
    %5th level
    \node[dotnode] (1234) at (12mm, -18mm){};
    \draw (123)--(1234)--(4);
    %root
    \node (root) at (12mm, -24mm){};
    \draw (1234)--(root);
  \end{tikzpicture}
  \qquad
  %%%%% ((12)(3 4)
  \begin{tikzpicture}
    %top level
    \node (1) at (0,0){};
    \node (2) at (8mm,0){};
    \node (3) at (16mm,0){};
    \node (4) at (24mm,0){};
    % labels
    \node[above, draw = none] at (1) {$a_1$};
    \node[above, draw = none] at (2) {$a_2$};
    \node[above, draw = none] at (3) {$a_3$};
    \node[above, draw = none] at (4) {$a_4$};
    % 2nd level
    \node[dotnode] (12) at (4mm, -6mm){};
    \draw (1)--(12)--(2);
    \node[dotnode] (34) at (20mm, -6mm){};
    \draw (3)--(34)--(4);
    %3rd level
    %4th level
    \node[dotnode] (1234) at (12mm, -18mm){};
    \draw (12)--(1234)--(34);
    %root
    \node (root) at (12mm, -24mm){};
    \draw (1234)--(root);
  \end{tikzpicture}
\end{center}
Adding ones requires a minor modification, adding appropriately
colored leaves to the tree.

For categories, the analogous structure is that of a monoidal category: a
category with a functor $\otimes \colon \calC\times \calC\to \calC$ and a unit
object $\one\in \calC$.

One can try to mimic the definition above.
\begin{definition}\label{d:strict_monoidal}
  A \emph{strict monoidal category}\index{monoidal category!strict} is a
  category $\calC$ together with a functor $\otimes \colon \calC\times \calC\to
  \calC$ and a unit object $\one\in \calC$ such that
  \begin{align*}
    &(A\otimes B)\otimes  C = A\otimes  (B\otimes C),\\
    &\one \otimes A  = A\otimes \one   = A.
  \end{align*}
\end{definition}

However, this definition is too strict: objects in a category are almost never
equal, just isomorphic. For example, if $V$ is a vector space over $\kk$, then
$\kk\otimes V\ne V$ (but is canonically isomorphic to $V$).

We can try to replace equality by isomorphism, but then it is too weak: e.g.,
any two vector spaces of the same dimension are isomorphic.

The proper way is to require the existence of a distinguished (canonical)
isomorphism which should be made part of the structure. Moreover, these
``canonical'' isomorphisms should themselves satisfy some compatibility
conditions, as suggested by Convention~\ref{conv:canonical}.

\begin{definition}\label{d:monoidal}\index{monoidal category}
  A \emph{monoidal category} is a category $\calC$ together with the following
  data:
  \begin{itemize}
    \item an object $\one \in \calC$
    \item a functor $\otimes\colon \calC \times \calC \to \calC$
    \item functorial isomorphisms
      \begin{align*}
        &l_A\colon \one \otimes A \isoto A\\
        & r_A\colon A\otimes \one \isoto A\\
        &\al_{A,B,C}\colon (A\otimes B)\otimes C\isoto A\otimes (B\otimes C)
      \end{align*}
  \end{itemize}
  satisfying the following compatibility conditions:
  \begin{enumerate}
    \item Triangle axiom: for any $A,B\in \calC$, the diagram below commutes
      \begin{equation}\label{e:triangle}
        \begin{tikzcd}[column sep=5mm, row sep=large]
          (A\otimes \one)\otimes B
          \arrow[rr,"\alpha_{A,\one,B}"]
          \arrow[dr,"r_A\otimes \mathrm{id}_B"']
          && A\otimes(\one\otimes B)
          \arrow[dl,"\mathrm{id}_A\otimes l_B"] \\
          & A\otimes B &
        \end{tikzcd}
      \end{equation}
    \item Pentagon axiom: for any $A,B,C,D \in \calC$, the diagram below
      commutes
      \begin{equation}\label{e:pentagon}
        \begin{tikzpicture}[vcenter]
          \node[scale=0.8] (a1) at (-3.5,0)
                      {$((A\otimes B)\otimes C)\otimes D$};
          \node[scale=0.8] (a2) at (-2,1.5)
                      {$(A\otimes (B\otimes C))\otimes D$};
          \node[scale=0.8] (a3) at (2,1.5)
                      {$A\otimes ((B\otimes C)\otimes D)$};
          \node[scale=0.8] (a4) at (3.5,0)
                      {$A \otimes (B \otimes (C \otimes D))$};
          \node[scale=0.8] (a5) at (0,-1.3)
                      {$(A \otimes B) \otimes (C \otimes D)$};
          \draw[->] (a1)--(a2);
          \draw[->] (a2)--(a3);
          \draw[->] (a3)--(a4);
          \draw[->] (a1)--(a5);
          \draw[->] (a5)--(a4);
        \end{tikzpicture}
      \end{equation}
      (all arrows are given by the associativity isomorphism $\al$, tensored
      with identity morphisms where necessary.)
%      \begin{equation}\label{e:pentagon}
%        \begin{tikzcd}[column sep=huge,row sep=large]
%          ((A\otimes B)\otimes C)\otimes D
%          \arrow[r,"\alpha_{A,B,C}\otimes \mathrm{id}_D"]
%          \arrow[d,"\alpha_{A\otimes B,C,D}"']
%          & (A\otimes (B\otimes C))\otimes D
%          \arrow[r,"\alpha_{A,B\otimes C,D}"]
%          & A\otimes((B\otimes C)\otimes D)
%          \arrow[d,"\mathrm{id}_A\otimes \alpha_{B,C,D}"] \\
%          (A\otimes B)\otimes(C\otimes D)
%          \arrow[rr,"\alpha_{A,B,C\otimes D}"']
%          && A\otimes(B\otimes(C\otimes D))
%        \end{tikzcd}
%      \end{equation}

  \end{enumerate}
\end{definition}

\begin{example}\label{x:monoidal-cat-examples}
The following are examples of monoidal categories.
  \begin{enumerate}
    \item Category $\sets$, with $\otimes$ given by disjoint union and
      $\one$ being the empty set
    \item Category $\sets$, with $\otimes$ given by the Cartesian product
      and $\one$ being a one-point set
    \item Categories $\smoothman$, $\plman$ with respect to disjoint union
    \item Category $\vect$ with the usual tensor product, $\one = \kk$
    \item Category of bimodules over a fixed ring $R$, with respect
      to $\otimes_R$, and $\one = R$

  \end{enumerate}
\end{example}

\begin{example}\label{x:fund-group}
  Let $G$ be a topological group. Let $\Pi_1(G)$ be its fundamental groupoid,
  as defined in \exref{x:fundamental-groupoid}. Then multiplication on $G$
  defines on $\Pi_1(G)$ a structure of a monoidal category.
\end{example}


\begin{example}\label{x:vecg}
  Let $G$ be a group. Define the category $\vect_G$ of \emph{$G$-graded vector
  spaces}\index{G-graded vector spaces@$G$-graded vector spaces} as the category
  whose objects are finite-dimensional vector spaces $V$ together with
  a decomposition
  $$
    V=\bigoplus_{g\in G} V_g.
  $$
  This is a monoidal category with tensor product given by
  $$
    (V\otimes W)_g = \bigoplus_{h\in G} V_h \otimes W_{h^{-1}g}
  $$
  and unit object given by $\kk$ considered as a $G$-graded vector space with
  grading equal to $1\in G$. The associativity isomorphism is trivial.
\end{example}



\begin{theorem}[Mac Lane's coherence theorem]\label{t:maclane}\par

  \noindent Let $\calC$ be a monoidal category. Let
  $F_1, F_2\colon \calC^n\to \calC$ be two functors obtained by compositions of
  the tensor product functor $\otimes$ and functors $X\mapsto X\otimes \one$,
  $X\mapsto \one \otimes X$, e.g.,
  \begin{align*}
    F_1(A,B,C)&=(A\otimes B)\otimes (\one\otimes C),\\
    F_2(A,B,C)&=A\otimes (( B\otimes  C)\otimes \one).
  \end{align*}
  Then
  \begin{enumerate}
    \item There exists a functorial isomorphism $F_1\simeq F_2$, obtained
      by composition of the associativity isomorphism $\al$, the left and right
      unit isomorphisms $l,r$, and their inverses.
    \item Any two such isomorphisms $F_1\to F_2$ are equal.
  \end{enumerate}
\end{theorem}
In other words, if we construct a cell complex $X_n$ whose vertices are functors
$\calC^n\to \calC$ as in the statement of the theorem, edges correspond to
applications of the associativity and unit isomorphisms, and 2-cells correspond
to commutative diagrams given by the pentagon axiom, unit axioms, and
functoriality of $\al, l, r$, then this complex is simply connected.

This is a natural generalization of the coherence theorem for monoids
(\thref{t:coherence-monoids}); in this language, \thref{t:coherence-monoids}
is exactly claim (1) of the theorem, i.e., the statement that $X_n$
is connected.

\begin{remark}\label{r:stasheff}
  If we only consider functors obtained by composing the tensor product functor
  (thus, no units), the corresponding complex is the 2-skeleton of a convex
  polytope $K_n$, called the \emph{associahedron}\index{associahedron}, or the
  \emph{Stasheff
  polytope}\index{Stasheff polytope}. The vertices of $K_n$ are indexed by
  binary trees with $n$ leaves; thus, the number of vertices is the Catalan
  number $C_{n-1} = \frac{1}{n}\binom{2(n-1)}{n-1}$ (so $K_5$ has
  $C_4 = 14$ vertices).
  This polytope was introduced by Jim Stasheff in the 1960s. It can be realized
  as a convex polytope in Euclidean space; see, e.g.,
  \ocite{loday}.
\end{remark}

\begin{exercise}\label{ec:K5}
  Draw the associahedron $K_5$ (it has 14 vertices) and verify that its
  2-skeleton is simply connected (in fact, it is homeomorphic to the sphere
  $S^2$). We will return to it later, when talking about monoidal 2-categories.
\end{exercise}

\begin{definition}\label{d:monoidal_functor}\index{monoidal functor}
  Let $\calC, \calD$ be monoidal categories. A monoidal functor
  $\calC\to \calD$ is a functor $F\colon\calC\to\calD$ together
  with functorial isomorphisms
  \begin{align*}
    J&\colon F(X\otimes Y)\to F(X)\otimes F(Y)\\
    \ph&\colon \one_\calD\to F(\one_\calC)
  \end{align*}
  satisfying suitable constraints (see \ocite{etingof-fusion}*{Section 2.4}).

  We say that categories $\calC, \calD$ are monoidally equivalent
  if there exist monoidal functors $F\colon \calC\to \calD$, $G\colon \calD\to
  \calC$ such that $F\circ G\simeq \id_\calD$, $G\circ F \simeq \id_\calC$.
\end{definition}

\begin{remark}\label{r:lax}
  In many references, monoidal functors as defined above are called \emph{weak
  monoidal functors}. One can also consider variations of this definition:
  instead of requiring an isomorphism $F(X\otimes Y)\simeq F(X)\otimes F(Y)$,
  one can require a functorial morphism (not necessarily invertible)
  $F(X\otimes Y)\to F(X)\otimes F(Y)$, or a functorial morphism in the opposite
  direction.  This gives the definitions of a \emph{colax/lax monoidal functor}
  respectively; see \ocite{leinster}*{Definition 1.2.10}.
\end{remark}


\begin{example}\label{x:rep-group-monoidal}
  Consider the category of representations of a group $G$. Then the
  forgetful functor from this category to the category of vector
  spaces is a monoidal functor. More generally, given
  a subgroup $H\subset G$, the restriction functor
  $\Res^G_H\colon \Rep G\to\Rep H$ has an
  obvious monoidal structure.
\end{example}
\begin{example}\label{x:rep-lie-group}
  If $G$ is a connected, simply-connected Lie group, then the category
  of finite-dimensional representations of $G$ is equivalent, as a
  monoidal category, to the category of finite-dimensional
  representations of its Lie
  algebra $\g$.
\end{example}

\begin{example}\label{x:vecg-twisted}
  Recall the category of $G$-graded vector spaces $\vect_G$ defined in
  \exref{x:vecg}. This category has a generalization. Namely, let
  $\om\colon G\times G\times G\to \kk^\times$ be a 3-cocycle, i.e., a
  function satisfying the condition
  $$
  \om(g_1g_2, g_3, g_4)\om (g_1, g_2, g_3g_4) = \om(g_1, g_2, g_3)
     \om(g_1, g_2g_3, g_4)\om (g_2, g_3, g_4).
  $$
  Consider the category $\vect^\om_G$ which coincides with $\vect_G$ as an
  abelian category; moreover, it has the same tensor product functor $\otimes$
  and unit object $\one=\kk_1$. However, the associativity and unit isomorphisms
  are different: namely, $\al \colon (V\otimes W)\otimes U \to V \otimes
  (W\otimes U)$ is given by
  $$
    \al((v \otimes w) \otimes u )= \om (g,h,m) v \otimes (w\otimes u),
    \qquad v\in V_g, w\in W_h, u\in U_m,
  $$
  and the unit isomorphisms $l_V\colon \one\otimes V\to V$,
  $r_V\colon V\otimes \one\to V$ are given by
  $$
    l_V(1\otimes v)=\om(1,1,g)^{-1}\, v, \qquad r_V(v\otimes 1)=\om(g,1,1)\, v,
    \qquad v\in V_g.
  $$
  The cocycle condition immediately implies that the $\al$ so defined satisfies
  the pentagon axiom. Moreover, taking $(g_1,g_2,g_3,g_4)=(g,1,1,h)$ in the
  cocycle condition gives $\om(g,1,h)=\om(g,1,1)\,\om(1,1,h)$, which is exactly
  the triangle axiom \eqref{e:triangle} for these $\al$, $l$, $r$. (If
  $\om(g,1,h)=1$ for all $g,h$, then $l$ and $r$ are the usual unit isomorphisms
  of $\vect_G$; for a general 3-cocycle, however, the usual unit isomorphisms
  would violate the triangle axiom.)
\end{example}
\begin{exercise}\label{ec:twisted-group-cat}
  Let $\om\colon G\times G\times G\to \kk^\times$ be a 3-cocycle. Assume that
  $\om$ is a coboundary, i.e., there exists a function $\be\colon G\times
  G\to \kk^\times$ such that $\om=d\be$, where the coboundary operator $d$
  is given by
  $$
    (d\be)(g_1, g_2, g_3)=
    \frac{\be(g_2, g_3)\,\be(g_1, g_2g_3)}{\be(g_1g_2, g_3)\,\be(g_1, g_2)}.
  $$
  Show that then the category $\vect_G^\om$ is monoidally equivalent to
  $\vect_G$.

  More generally, the categories $\vect_G^\om$ and $\vect_G^{\om'}$ are
  monoidally equivalent, with the equivalence preserving the $G$-graded
  dimension, if and only if $\om'=\om \cdot d\be$ for some
  $\be\colon G\times G\to \kk^\times$ (see \ocite{etingof-fusion}*{Section
  2.6}).
\end{exercise}



We can restate Mac Lane's coherence theorem as follows.
\begin{theorem}\label{t:strictify}
  Any monoidal category is monoidally equivalent to a strict monoidal category.
\end{theorem}




\begin{convention}\label{c:graphical-monoidal}
  We will frequently use the graphical presentation of morphisms
  in a monoidal category, representing a morphism
  $\ph \colon A_1\otimes \dots \otimes A_k\to B_1\otimes\dots \otimes B_m$
  by a box with $k$ inputs at the top and $m$ outputs at the bottom as shown in
  \firef{f:graphical-morphism}.
\begin{figure}[h]
  \centering
  \begin{tikzpicture}
    \node[morphism_sq, minimum width=2cm] (ph) at (0,0) {$\ph$};
    %% top
    \node[label=above:$A_1$] (A1) at (-0.7,1){};
    \node[label=above:$A_2$]  (A2) at (0,1){};
    \node[label=above:$A_3$]  (A3) at (0.7,1){};
    \draw (A1)-- (A1|-ph.north);
    \draw (A2)-- (A2|-ph.north);
    \draw (A3)-- (A3|-ph.north);
    %% bottom
    \node[label=below:$B_1$] (B1) at (-0.4,-1){};
    \node[label=below:$B_2$]  (B2) at (0.4,-1){};
    \draw (B1)-- (B1|-ph.south);
    \draw (B2)-- (B2|-ph.south);
    %% eqn
    \node[right] at (2,0) {$\ph\colon A_1\otimes A_2\otimes A_3\to B_1\otimes
    B_2$};
  \end{tikzpicture}
  \caption{Graphical notation for a morphism in a monoidal category.}
  \label{f:graphical-morphism}
\end{figure}
\end{convention}







%%%%%%%%%%%%%%%%%%%%%%%%%%%%%%%%%%%%%%%
\section{Symmetric monoidal categories}\label{s:symmetric-monoidal}
For the category of vector spaces, the tensor product is not only associative,
but also ``commutative'': one has natural isomorphisms $V\otimes W \simeq
W\otimes V$. The definition below gives a proper generalization of this to
arbitrary monoidal categories.

\begin{definition}\label{d:symmetric}\index{monoidal category!symmetric}
  A \emph{symmetric monoidal category} is a monoidal category $\calC$ together
  with a functorial isomorphism
  $$
    c_{X,Y}\colon X\otimes Y \isoto Y\otimes X
  $$
  such that
  $$
  c_{X,Y} c_{Y,X}=\id
  $$
  and which satisfies the hexagon axiom:
  \begin{equation}\label{e:hexagon1}
    \begin{tikzcd}[column sep=huge, row sep=large]
      (A\otimes B)\otimes C
      \arrow[r,"\alpha_{A,B,C}"]
      \arrow[d,"c_{A,B}\otimes \mathrm{id}_C"']
      & A\otimes(B\otimes C)
      \arrow[r,"c_{A,B\otimes C}"]
      & (B\otimes C)\otimes A
      \arrow[d,"\alpha_{B,C,A}"] \\
      (B\otimes A)\otimes C
      \arrow[r,"\alpha_{B,A,C}"']
      & B\otimes(A\otimes C)
      \arrow[r,"\mathrm{id}_B\otimes c_{A,C}"']
      & B\otimes(C\otimes A)
    \end{tikzcd}
  \end{equation}
\end{definition}
The hexagon axiom looks complicated; to make it easier to understand, one
can drop all associativity isomorphisms, replacing the hexagon axiom with the
following diagram:
\begin{equation}\label{e:hexagon3}
  \begin{tikzcd}[column sep=5mm, row sep=large]
    A\otimes B\otimes C
    \arrow[rr,"c_{A, B\otimes C}"]
    \arrow[dr,"c_{A,B}"']
    && B\otimes C\otimes A\\
    & B\otimes A\otimes C
    \arrow[ur,"c_{A,C}"]  &
  \end{tikzcd}
\end{equation}

\begin{example}\label{x:symmetric-monoidal-cats}\par\noindent
  \begin{enumerate}
    \item Category $\sets$ with respect to disjoint union is a
      symmetric monoidal category; the same is true for the categories
      $\smoothman$, $\plman$.
    \item
      The category of vector spaces over $\kk$ is a symmetric monoidal category.
      The same is true for the categories of representations of a group $G$ or
      a Lie algebra $\g$.
    \item (Super vector spaces). Consider the category $\svect$ whose
      objects are $\Z_2$-graded vector spaces $V=V_0\oplus V_1$. This
      has an obvious monoidal structure.
      Define the commutativity isomorphism $c_{V,W}\colon V\otimes W
      \to W\otimes V$ by
      $$
      c(v\otimes w)=(-1)^{pq} w\otimes v, \quad v\in V_p, w\in W_q.
      $$
      Then this defines on the category $\svect$ the structure of a
      symmetric monoidal category.
  \end{enumerate}
\end{example}
Note that one can also define a symmetric structure on $\svect$ without the
$(-1)^{pq}$ factor; this illustrates the fact that a monoidal category can have
several different symmetric structures.
\begin{remark}\label{r:non-symmetric}
  There are monoidal categories that do not have a natural symmetric structure;
  for example, the category $\vect_G$ of $G$-graded vector spaces for a
  non-abelian group $G$, or a category of bimodules over a ring. There are also
  ``braided'' categories where we do have isomorphisms
  $c_{X,Y}\colon X\otimes Y \to Y\otimes X$, but these isomorphisms fail to
  satisfy the condition $c^2=\id$. The most famous example of a braided category
  is the category of representations of a quantum group. We will discuss braided
  categories later (see \chref{c:321}).
\end{remark}

As before, the real importance of the hexagon axiom is the following coherence
theorem, which is an analog of Mac Lane's coherence theorem, but for
\textbf{symmetric} monoidal categories.

\begin{theorem}[Mac Lane's coherence theorem for symmetric monoidal
 categories]\label{t:maclane-symmetric}\par

  \noindent
  Let $\calC$ be a symmetric monoidal category. Let $F_1, F_2\colon \calC^n\to
  \calC$ be two functors obtained by compositions of
  \begin{itemize}
    \item the tensor product functor $\otimes\colon \calC\times \calC\to \calC$
    \item the functors $X\mapsto X\otimes \one$, $X\mapsto \one \otimes X$
    \item the permutation functor $\calC\times \calC \to \calC\times \calC
          \colon (A, B)\mapsto (B,A)$
  \end{itemize}
  e.g.,
  \begin{align*}
    F_1(A,B,C)&=(B\otimes A)\otimes (\one\otimes C)\\
    F_2(A,B,C)&=C\otimes ((B\otimes  A)\otimes \one).
  \end{align*}
  Then
  \begin{enumerate}
    \item There exists a functorial isomorphism $F_1\simeq F_2$, obtained
      by composition of the associativity isomorphism $\al$, the left and right
      unit isomorphisms $l,r$, the commutativity isomorphism $c$, and their
      inverses.
    \item Any two such isomorphisms $F_1\to F_2$ are equal.
  \end{enumerate}
\end{theorem}
Similar to \thref{t:maclane}, this theorem can be understood as a statement
that a certain cell complex is simply connected; if we only use
permutations and tensor product (i.e., no units), then this complex is
called the \emph{permutoassociahedron}. (Yes, we know, it is long and
impossible to spell.)


As before, we skip the proof of this theorem, referring the reader to
\ocite{mac-lane}*{Section XI.1}.

\begin{corollary}\label{c:AS}
  Let $\calC$ be a symmetric monoidal category and let $A_s$, $s\in S$, be a
  collection of objects in $\calC$ indexed by a finite set $S$. Then we have a
  well-defined (up to a canonical isomorphism) object $ \bigotimes_{s\in S}
  A_s$.
\end{corollary}
%%%%%%%%%%%%%%%%%%%%%%%%%
\section{Rigid monoidal categories}\label{s:rigid}
In many monoidal categories, such as the category of finite-dimensional vector
spaces or the category of finite-dimensional representations of a group $G$, in
addition to the tensor product we also have the operation of taking duals: for
every representation $V$, we also have the dual representation $V^*$. The
following definition allows one to talk about duals in any monoidal category.
\begin{definition}\label{d:dual-object}
  Let $\calC$ be a monoidal category.
  A \emph{dual pair}\index{dual pair} is a pair of objects $A,B\in \calC$
  together with morphisms
  \begin{align*}
    \ev&\colon A\otimes B\to \one \qquad \text{(evaluation morphism),}\\
    \coev&\colon \one \to B\otimes A \qquad \text{(coevaluation morphism)},
  \end{align*}
  satisfying the \emph{rigidity}\index{rigidity} condition below.
  \begin{equation}\label{e:rigidity}
    \begin{aligned}
      &\Bigl(B\to \one\otimes B \xto{\coev\otimes 1_B}
                       (B\otimes A)\otimes B \xto{\al_{B,A,B}}
                       B\otimes (A\otimes B) \xto{1_B\otimes \ev}
                       B\otimes \one \to
                       B\Bigr) = 1_B\\
      &\Bigl(A\to  A\otimes \one  \xto{1_A\otimes \coev}
                       A\otimes (B\otimes A) \xto{\al^{-1}_{A,B,A}}
                       (A\otimes B)\otimes A \xto{\ev\otimes 1_A}
                       \one\otimes A \to
                       A\Bigr) = 1_A
    \end{aligned}
  \end{equation}
  In this situation we say that $A$ is a
  \emph{left dual}\index{dual object!left} of $B$, and
  $B$ is a \emph{right dual}\index{dual object!right} of $A$.
\end{definition}
Note that in this definition we do not assume that $\calC$ is symmetric.


It is common to represent the evaluation and coevaluation graphically:
\begin{equation}\label{e:ev-coev-graphical}
  \begin{tikzpicture}
  \node at (0.3,0) {$\ev =$};
  \node[label=above:$A$] (A) at (1,0) {};
  \node[label=above:$B$] (B) at (1.5,0) {};
  \pic at (A) {ev};
  \node at (3.3,0) {$\coev =$};
  \node[label=below:$B$] (B) at (4,0) {};
  \node[label=below:$A$] (A) at (4.5,0) {};
  \pic at (B) {coev};
  \end{tikzpicture}
\end{equation}

Then the rigidity equations \eqref{e:rigidity} are illustrated by
\begin{equation}\label{e:rigidity2}
  \begin{tikzpicture}[vcenter]
    \node[label=above:$A$] (A) at (0,0){};
    \pic at (A) {rigidity};
    \node at (1.5, -1) {$=$};
    \node[label=above:$A$] (A2) at (2,0){};
    \draw (A2)-- +(0,-2);
  \end{tikzpicture}
  \qquad\qquad
  \begin{tikzpicture}[vcenter]
    \node[label=above:$B$] (B) at (0,0){};
    \pic[x=-1cm] at (B) {rigidity};
    \node at (0.5, -1) {$=$};
    \node[label=above:$B$] (B2) at (1,0){};
    \draw (B2)-- +(0,-2);
  \end{tikzpicture}
\end{equation}

The following example provides the motivation for this definition.

\begin{example}\label{x:dual-vector}
  Let $V$ be a finite-dimensional vector space and let $V^*$ be the dual vector
  space. Let $\ev\colon V^*\otimes V\to \kk$ be the natural pairing. Define
  the coevaluation
  \begin{equation}\label{e:dual-vector}
      \begin{aligned}
      \coev\colon \kk &\to V\otimes V^*\\
                  1&\mapsto \sum v_i\otimes v^i
      \end{aligned}
  \end{equation}
  where $v_i, v^i$ are dual bases in $V, V^*$ respectively. Then the maps $\ev$
  and $\coev$ so defined satisfy the rigidity equations \eqref{e:rigidity};
  thus, $(V^*, V, \ev, \coev)$ form a dual pair.

  Indeed, one easily sees that the rigidity equations are equivalent to
  requiring that for any $v\in V, f\in V^*$, we have
  \begin{equation}
   \begin{aligned}
    \sum \<v^i, v\>v_i &= v\\
    \sum \<f, v_i\>v^i &= f
   \end{aligned}
  \end{equation}
\end{example}

\begin{lemma}\label{l:rigidity_vect}
  A vector space $V$, considered as an object in the monoidal category $\vect$,
  has a left dual if and only if it is finite-dimensional. In this case, the
  left dual is given by $V^*$, with evaluation and coevaluation defined as in
  \exref{x:dual-vector}.
\end{lemma}

We now return to dual objects in arbitrary monoidal categories.

It turns out that the right dual, if it exists, is unique up to a unique
isomorphism.
\begin{theorem}\label{t:dual_unique1}
  Let $(A, B, \ev, \coev)$ and $(A,B', \ev', \coev')$ be two dual pairs with
  the same $A$. Then there exists a unique morphism $\ph\colon B\to B'$ such
  that $\ev = \ev' \circ (1_A\otimes \ph)$, $\coev' =(\ph\otimes 1_A) \coev$.
  Moreover, $\ph$ is invertible. 
\end{theorem}
The proof is a standard snake-diagram argument: the required $\ph$ is given by
the composition
$$
  B \xto{\coev'\otimes 1_B} B'\otimes A\otimes B
    \xto{1_{B'}\otimes \ev} B',
$$
and uniqueness follows by chasing the rigidity equations.
A similar statement holds for left duals.

This allows us to talk about the right dual of $A$ as a well-defined object in
$\calC$. We will denote the right dual of $A$ by ${}^*A$; similarly we will use
the notation $B^*$ for the left dual of $B$. In this notation, the evaluation
and coevaluation morphisms become
\begin{equation}
  \begin{aligned}
    \ev_B&\colon B^*\otimes B\to\one\\
    \coev_B&\colon \one \to B\otimes B^*.
  \end{aligned}
\end{equation}

Clearly, in a symmetric monoidal
category, the left dual is also a right dual.

\begin{definition}\label{d:rigid}
  An object $A$ in a monoidal category $\calC$ is
  \emph{rigid}\index{rigid object} if
  it has both left and right duals.

  A monoidal category is \emph{rigid}\index{monoidal category!rigid}
  if every object is rigid.
\end{definition}

\begin{example}\label{x:vector-rigid}
  By \leref{l:rigidity_vect}, a vector space $V$ is rigid as an object of
  $\vect$ if and only if it is finite-dimensional.
\end{example}


\begin{exercise}\label{ec:coev-unique}
  Show that in a dual pair, the coevaluation is uniquely determined by
  the remaining data: if $(A, B, \ev, \coev)$ and $(A,B, \ev, \coev')$
  are dual pairs, then $\coev = \coev'$.

  Similarly, $\ev$ is uniquely determined by $(A,B, \coev)$.
\end{exercise}

\begin{exercise}\label{ec:dual-of-tensor}
  Let $X^*$, $Y^*$ be left duals of objects $X, Y$. Show that then $X\otimes Y$
  also has a left dual, and there is a canonical isomorphism $(X\otimes
  Y)^*\simeq Y^* \otimes X^*$.
\end{exercise}

\begin{exercise}\label{ec:dual-morphism}
  Let $\calC$ be a rigid monoidal category. For a morphism $f\colon X\to Y$,
  define a morphism $f^*\colon Y^*\to X^*$ by the picture below.

  \begin{equation}\label{e:dual-morphism}
    \begin{tikzpicture}[vcenter]
      \node at (-1.2,0) {$f^*=$};
      \node[morphism_sq] (f) at (0,0) {$f$};
      \draw(f.north)-- +(0, 0.25) arc (180:0:0.25) -- +(0, -1.5) coordinate
      (X*);
      \draw(f.south)-- +(0, -0.25) arc (0:-180:0.25) -- +(0, 1.5) coordinate
      (Y*);
      \node[above] at (Y*) {$Y^*$};
      \node[below] at (X*) {$X^*$};
    \end{tikzpicture}
  \end{equation}


  Prove that then
  $(fg)^*=g^* f^*$; thus, $*$ is a contravariant functor $\calC\to \calC$.
\end{exercise}


\begin{exercise}\label{ec:dual-morphism2}
  Let $\calC$ be a rigid monoidal category. Show that then one has
  canonical functorial isomorphisms
  \begin{align*}
    \Hom(X,Y)&\simeq \Hom(\one, Y\otimes X^*)\simeq \Hom(\one, {}^*X\otimes Y)\\
    \Hom(X\otimes Z, Y)&\simeq \Hom(Z, {}^*X\otimes Y)\simeq \Hom(X, Y\otimes
    Z^*).
  \end{align*}
\end{exercise}

